\subsection{Semipositive residuals with positive volume}
\label{sec:semipositive-fiber-integrability}

The divisor-order estimates do not use strict positivity of the smooth residual. Positivity enters only when constructing the reference current and ensuring that its non-pluripolar mass is nonzero. This separates the section comparison from the existence of a K\"ahler representative of the residual class.

\Needspace{12\baselineskip}
\begin{corollary}\label{cor:semipositive-fiber-band}\index{Holomorphic section!semipositive residual}
In each of \cref{thm:fiber-integrability-band,thm:crossing-fiber-integrability-band,thm:modified-base-fiber-band}, replace the condition $\beta>0$ by
\[
 \beta\geq0,\qquad \int_M\beta^q>0,
\]
where $M=B$ in the first two theorems and $M=Y$ in the third. Retain every other hypothesis. All the conclusions remain valid. The same replacement is valid in \cref{cor:exact-anisotropic-fiber-series}.
\end{corollary}
\begin{proof}
In each case the reference current is $\beta+\sum_i c_i[D_i]$, hence positive. Locality off the divisor union and exclusion of that union give the non-pluripolar measure $\beta^q$ on $M$. Smoothness and compactness make its mass finite, and the assumed top integral makes it positive. In the modified-base case the potential descends and its measure is $\rho_*(\beta^q)$ by the same descent and measure identities.

The base multiplier ideal depends only on the logarithmic divisor orders and the bounded smooth term. The lower inclusion uses the global selector bound; the upper inclusion uses the anisotropic threshold in a witness chart. Neither requires strict positivity of $\beta$. The restriction-dimension estimates apply to fixed line bundles regardless of their curvature. The Jacobian and fixed canonical twist in the modified-base argument are unchanged. Thus the inclusions and dimension bounds in all three proofs still hold. The reverse inclusion for the exact anisotropic model uses only its local integrability test and local boundedness off the pole set, so it also remains valid.
\end{proof}

The positive top integral is a condition on the chosen reference. Semipositivity alone permits zero mass. Conversely, positive top integral need not allow any smooth gauge to make the residual strictly positive. The following example realizes this distinction with an exact selected series.

\Needspace{14\baselineskip}
\begin{example}\label{ex:exceptional-semipositive-fiber-series}\index{Holomorphic section!exceptional fixed divisor}
Let $q\geq2$, let $\rho:B\to\mathbb P^q$ be the blowup of $p=[1:0:\cdots:0]$, and let $E$ be its exceptional divisor. Let $\tau:B\to\mathbb P^{q-1}$ resolve the projection from $p$. Put
\[
 \alpha=\rho^*\omega_{\mathbb P^q},\qquad
 \gamma=\tau^*\omega_{\mathbb P^{q-1}},\qquad
 \kappa=\alpha+\gamma>0.
\]
Use the quotient metric on $\mathcal O_B(E)=\rho^*\mathcal O_{\mathbb P^q}(1)\otimes\tau^*\mathcal O_{\mathbb P^{q-1}}(-1)$, with curvature $\eta=\alpha-\gamma$ and defining section $s_E$. For an integer $k\geq1$, equip
\[
 A=\rho^*\mathcal O_{\mathbb P^q}(k)\otimes\mathcal O_B(E)
 \quad\hbox{with curvature}\quad \theta=k\alpha+\eta.
\]
Take $X=B\times\mathbb P^1$, with projection $\pi$ to $B$. For a marked point $o$ in the fiber, use its normalized Fubini--Study divisor metric and set
\[
 \begin{gathered}
 R=\|\pi^*s_E\|^2,\quad T=\|s_o\|^2,\quad
 \varphi=\tfrac12\log(R^2+T^3),\quad
 \omega=\pi^*\kappa+2\omega_{\mathbb P^1},\\
 \psi=\log\|s_E\|^2.
 \end{gathered}
\]
Then $\varphi\in\PSH(X,\omega)$, $\psi\in\PSH(B,\theta)$ and
\[
 \mathrm{MA}_\theta(\psi)=(k\alpha)^q,\qquad
 \int_B\mathrm{MA}_\theta(\psi)=k^q.
\]
The residual $k\alpha$ admits no strictly positive representative obtained by a smooth gauge. Nevertheless, for every $m\geq1$ the selected and reference spaces coincide as subspaces of $H^0(B,A^m)$:
\[
 V_m=W_m=H^0(B,A^m)
      =s_E^m\rho^*H^0(\mathbb P^q,\mathcal O(km)),
 \qquad \dim V_m=\binom{km+q}{q}.
\]
\end{example}
\begin{proof}
The graph embedding $(\rho,\tau)$ gives $\kappa>0$. The quotient metric satisfies
\[
 \|s_E\|^2=\rho^*\left(\frac{\sum_{j=1}^q|Z_j|^2}
                                  {\sum_{j=0}^q|Z_j|^2}\right),
\]
as in \cref{ex:modified-point-fiber-series}. The same common-denominator calculation, now with coefficient $1/2$, gives
\[
 \ddc\varphi\geq-\pi^*\alpha-\tfrac32\omega_{\mathbb P^1}.
\]
After local psh extension across $E\times\{o\}$, this implies $\omega+\ddc\varphi\geq\pi^*\gamma+\tfrac12\omega_{\mathbb P^1}\geq0$. The selector has the exact model $\tfrac12\log(|z|^4+|w|^6)+O(1)$ along that smooth pole set and is locally bounded elsewhere. It also satisfies the global lower bound $\varphi\geq\log R$.

Here $c=1$, $u=0$, $p=2$, $b_1=3$, $a=2/3$ and $\beta=\theta-\eta=k\alpha\geq0$. Its top integral is $k^q$, since the normalized Fubini--Study top integral is one and $\rho$ is a modification. \Cref{cor:semipositive-fiber-band,cor:exact-anisotropic-fiber-series} therefore give
\[
 W_m=H^0(B,A^m(-mE)),\qquad
 V_m=H^0(B,A^m(-(m-1)E)).
\]

For every integer $t\geq0$, a section of $\rho^*\mathcal O(km)\otimes\mathcal O_B(tE)$ divided by $s_E^t$ gives a holomorphic section of $\mathcal O(km)$ on $\mathbb P^q\setminus\{p\}$. Hartogs extension applies because $q\geq2$. Pulling the extension back and multiplying by $s_E^t$ recovers the original section off $E$, hence everywhere by the identity theorem. Thus all sections of $A^m$ are $s_E^m\rho^*s$. The natural embeddings of the two displayed spaces into $H^0(B,A^m)$ both have this image. Their dimension is the homogeneous-monomial count $\binom{km+q}{q}$.

Finally, a projective line $C\subset E\simeq\mathbb P^{q-1}$ is contracted by $\rho$. For every smooth real $u$, Stokes' theorem gives $\int_C(k\alpha+\ddc u)=0$, whereas a strictly positive form would have positive integral on $C$. Changing the smooth divisor metric only changes this residual by a smooth $\ddc$ term. This proves the claimed obstruction. The selector is finite on most of each fiber over $E$ and equals $-\infty$ at $o$, so it does not descend up to bounded difference.
\end{proof}

The example has zero section defect despite its non-descended selector and positive reference mass. Its fixed exceptional divisor forces all global sections to vanish more than the upper local threshold demands. The criterion still requires the global selector bounds and a specified positive-volume residual; it does not furnish them for arbitrary fiber singularities.
